\documentclass{article}
\usepackage[T1]{fontenc}
\usepackage{lmodern}
\usepackage{graphicx}
\usepackage{amsmath,amssymb}
\usepackage[a4paper,margin=2cm]{geometry}
\usepackage[absolute,overlay]{textpos}
\setlength{\TPHorizModule}{1mm}
\setlength{\TPVertModule}{1mm}
\usepackage[indonesian]{babel}
\usepackage{hyperref}
\setlength{\emergencystretch}{2em}
% unit: Halaman 90

\begin{document}

% === Halaman 90 (python/native) ===

90

• Affine Cipher:

             Enkripsi: C  mP + b (mod n) 


Dekripsi: P   m–1 (C – b) (mod n)

• Affine cipher mudah diserang dengan known-plaintext attack (sebagian plainteks dari 

pesan diketahui) 

• Misalkan kriptanalis mempunyai dua buah plainteks, P1 dan P2, yang berkoresponden 

dengan cipherteks C1 dan C2, 

• maka m dan b mudah dihitung dari buah kekongruenan simultan berikut ini:


C1  mP1 + b (mod n) 


C2  mP2 + b (mod n)  

Kriptanalisis Affine Cipher

\end{document}
